\section{Evaluation Setup}
\label{sec:setup}
\subsection{Inputs}
The base feature set of each link (left number in the last column of Table~\ref{tab:data}) holds the three targets, the iperf throughput, the link telemetry listed in Section~\ref{sec:system}, and for Starlink and 5G three weather variables. The extended set adds, for each target, seven causal statistics computed within the session: the number of rows since the value last changed (the phase in the report cycle, since jitter and loss are one-second reports held for five rows), the mean and maximum of the last second, the means of the last 30\,s and 60\,s, the standard deviation of the last 60\,s, and the session-to-date mean. It also adds the time since the previous row. Radio values outside physical ranges, such as the logger sentinel $-32768$ in some 5G SINR records, are treated as missing, and non-target inputs are clipped to $\pm 5$ standard deviations after normalization with training statistics.

\subsection{Forecasters}
\begin{itemize}
  \item \textbf{Persistence} $\hat{\mathbf{y}}_{t+h}=\mathbf{y}_t$ and \textbf{Seasonal}, the value one second before the target.
  \item \textbf{Ridge-E}: a linear model per target and step that predicts $\mathbf{y}_{t+h}-\mathbf{y}_t$ from the last extended input row and the 30-row target history, with the penalty chosen on the calibration sessions.
  \item \textbf{LQR}: a linear map from the Ridge-E design to the outputs of one DS2S head, anchored and ordered as in~\eqref{eq:anchor} and trained with the loss~\eqref{eq:loss}. It is a linear quantile regression~\cite{koenker1978regression} with non-crossing quantiles, the linear twin of DS2S (one seed).
  \item \textbf{S2S}: the GRU encoder-decoder of Section~\ref{sec:method} with one scalar head per target, a decoder started from zero, the base features, and the squared-error loss. This is the forecaster of our earlier deployment of the prediction-aware scheduler, retrained here on causal inputs.
  \item \textbf{S2S-AE}: S2S with the anchored decoder and the extended features, still with one scalar head. It differs from DS2S only in the output and the loss.
  \item \textbf{GS2S}: DS2S with Gaussian heads that output a mean and an SD per step and target, trained with the Gaussian negative log-likelihood~\cite{nix1994estimating}. Its quantiles are those of the normal distribution and are therefore symmetric.
  \item \textbf{DS2S}: the proposed forecaster. A suffix -E denotes the ensemble of three seeds.
\end{itemize}
All neural models use AdamW with learning rate and weight decay $10^{-3}$, batch size 512, gradient clipping at norm 5, and a plateau schedule that halves the learning rate after 8 epochs without improvement. The checkpoint with the lowest exponentially smoothed calibration-session loss (factor 0.3) is kept, and training stops after 15 epochs without improvement or at 60 epochs. All \NRuns{} runs of the main comparison stopped early, after \EpochsMin{} to \EpochsMax{} epochs. Each neural method is trained with three seeds on an Apple M3 Max GPU.

\subsection{Intervals}
Every 95\% interval is calibrated on the calibration sessions only. A point forecaster gets one of three intervals. The \emph{constant} interval $\hat{y}\pm c_h$ uses the split conformal quantile of $|y-\hat{y}|$~\cite{lei2018distribution}. The \emph{asymmetric} interval $[\hat{y}-c^{-}_h,\hat{y}+c^{+}_h]$ uses the conformal quantiles of the signed error with $\alpha/2$ per tail. The \emph{SD-scaled} interval $\hat{y}\pm c_h(\hat{\sigma}+\varepsilon)$ uses the quantile of $|y-\hat{y}|/(\hat{\sigma}+\varepsilon)$, where $\hat{\sigma}$ is the SD of the 30 input RTT values and the floor $\varepsilon$ is chosen on the calibration sessions by the interval score. It is the strongest volatility heuristic we could give a point forecast. A quantile forecaster (LQR, GS2S, DS2S) gets the CQR interval of Section~\ref{sec:method}.

\subsection{Metrics and Statistics}
For each target and step we compute the MAE (with $\hat{q}_{0.5}$ for quantile forecasters), the RMSE (with $\hat{\mu}$), the coverage (PICP) of the 95\% interval, and the interval score~\eqref{eq:is}, and we average over the 15 steps. To measure conditional coverage, we rank the test windows into deciles of the input SD $\hat{\sigma}$, with decile edges from the calibration sessions, and report the coverage per decile and step and the largest decile deviation $\max_d|\mathrm{cov}_d-0.95|$ (coverage gap). The test windows of one session overlap and share its regime. The unit of independent evidence is therefore the session, not the window. We compare models with a paired bootstrap that resamples test sessions with replacement (2000 replicates, percentile intervals) and report a change as significant when its 95\% interval excludes zero.
